% 图 9.1 —— 手工重建（两个方框 B / A−B，B 上点列 a₁…a₅ 配箭头，映射 g）
%
% ══ 这一版是**按用户意见重做**的，原稿是 `emit_hand.py` 的自动稿，问题成串 ══
%   ① `g` 被认成 `9`：左侧 (328.3,64.8) 与右侧 (924.5,83.4) 两处
%   ② 下标认错：a₄ 的下标 4 画成 6、a₅ 的下标 5 画成 9；a₂/a₃ 的下标干脆没有
%   ③ 字号不统一：a₂/a₃ 用了 51.3pt，其余 29~31pt（差 1.7 倍）
%   ④ 左右框的标签 `B` 被拆成两段曲线画出来（渲染成像 θ）
%   ⑤ 五支箭头是「三段曲线 + 七零八落的短线段拼出箭头头」→ 头是圆疙瘩
%
% 样式**一律对齐图 7.1**（用户指定）：
%   字母 Arial 粗斜体、下标 Arial 粗体（不斜），箭头 `-{Triangle[...]}`。
%   7.1 实测换算：字母 x 高 ÷ 0.519 = 字号；数字高 ÷ 0.716 = 字号。
%
% 参数来源：原书 600dpi 裁切 `src_hi/9.1.png`（1882×489），逐块连通域实测。
%   标定：方框竖边 px 14..21 / 1278..1284 / 1857..1865，
%         横边 px 17..22 / 397..402
%         ⇒ 画布x = (px-17.5)×0.5465 + 17   ；画布y = (py-19.5)×0.5465 + 15
%         两路独立标定互相吻合，圆点实测 x 109.1/204.7/306.4/402.8/472.0/532.1/570.9/592.7/855.6
%         —— 与原稿 dots 逐个相符，故圆点与方框不动。
%   标签（实测块 → 画布中心）：字母高 16.4 单位 ⇒ 30pt；下标高 15.3 ⇒ 22pt；
%         g 高 23.5（含降部 0.731em）⇒ 32pt；框标签 B 帽高 21.3 ⇒ 30pt。
%   箭头：五条弧的墨迹框实测 x 跨度 68~70、y 75.9..98.9（弧顶 ≈76），
%         箭头头占末端约 18 单位 ⇒ Triangle[width=16pt,length=18pt]
\documentclass[12pt]{standalone}
\usepackage{fontspec}
\setsansfont{Arial}
\newfontfamily\mfallback{DejaVu Sans}
\usepackage{tikz}
\usetikzlibrary{arrows.meta}
\begin{document}
\begin{tikzpicture}[x=1pt, y=-1pt, line width=4.0pt, line cap=butt, line join=miter]

  %% ════ 两个方框（原稿是几十段碎线拼的，这里直接用矩形）════
  \draw[line width=4.0pt] (15.0,13.6) rectangle (709.0,224.0);
  \draw[line width=4.0pt] (709.0,13.6) rectangle (1026.6,224.0);

  %% ════ 点列（实测位置，与原稿一致，不动）════
  \foreach \x in {108.8,204.6,306.1,403.0,472.0,532.1,570.6,592.6}
    \fill (\x,126.6) circle (6.2);
  \fill (855.6,126.6) circle (6.2);

  %% ════ 五支箭头（左框：aᵢ → aᵢ₊₁）════
  %% ！每一条**都是一段圆弧**（+ 末端一个三角箭头头），不是自由贝塞尔。
  %% 2026-09-25 重做：上一版拿三次贝塞尔去拟「陡升 + 长平顶 + 陡降」，
  %%   一条三次曲线表达不了这种形状 —— 最小二乘残差 12 个单位，
  %%   渲染出来就是「顶点偏右、整条斜着」，和原书的"平顶拱"完全不像。
  %% 正解：逐列提中心线 → 只取**弧段列**（列宽 < 6.5，排掉箭头头）→ **最小二乘圆拟合**。
  %%   五条的圆度残差 **0.62 ~ 1.42** ⇒ 确认是圆弧。参数逐条实测如下：
  %%     ① 心(152.7,124.4) r=43.91  -141.1°→-43.0°   ② 心(255.5,125.0) r=43.99  -140.9°→-42.9°
  %%     ③ 心(351.9,122.7) r=41.95  -142.9°→-41.3°   ④ 心(434.9,118.6) r=36.14  -142.2°→-37.9°
  %%     ⑤ 心(500.7,113.1) r=30.01  -144.6°→-29.3°
  %%   半径从 43.9 递减到 30.0 —— 原书**后两条更扁**，不是同一套复制五遍。
  %% ⚠ 相邻两条的墨迹在 x 上会重叠（第 4 条的尖伸到第 5 条起点左边），逐列取样时必须
  %%   「取离预测曲线最近的一段」，不能取最宽的一段。
  %% 箭头头（2026-09-25 按用户意见第三版）：
  %% ① **头的尾巴放在弧的端点**——原来是把**尖**放在弧端，于是头把弧的最后 18~22 单位
  %%    整段盖住（用户："箭头的效果和线条有覆盖"）。现在弧画到「头的尾」就停，
  %%    头正好接在弧端外面，两者不再重叠。
  %% ② **头调小**：width 16 → 实测原书头的实心列墨框是 13×16.4 画布单位，
  %%    length=18 时渲染出来正好 13×16.4（length=22 会高到 19.7，偏大）。
  %% ③ 头的**方向**比弧末切向**逆时针浅 15°**（原书实测头轴 35~45°，而弧末切向 47~61°）；
  %%    做法：弧提前 stop，再补一段**与头等长**的直线到尖端，头的方向就取自这段直线的方向。
  %%    该直线被头完全盖住（长度相等）。
  \draw[line width=3.5pt,-{Triangle[width=16pt,length=18pt]}]
    (152.7,124.4) ++(-141.1:43.91) arc[start angle=-141.1, end angle=-66.89, radius=43.91] -- (184.8,94.5);
  \draw[line width=3.5pt,-{Triangle[width=16pt,length=18pt]}]
    (255.5,125.0) ++(-140.9:43.99) arc[start angle=-140.9, end angle=-66.75, radius=43.99] -- (287.7,95.1);
  \draw[line width=3.5pt,-{Triangle[width=16pt,length=18pt]}]
    (351.9,122.7) ++(-142.9:41.95) arc[start angle=-142.9, end angle=-66.30, radius=41.95] -- (383.4,95.0);
  \draw[line width=3.5pt,-{Triangle[width=16pt,length=18pt]}]
    (434.9,118.6) ++(-142.2:36.14) arc[start angle=-142.2, end angle=-66.81, radius=36.14] -- (463.4,96.4);
  \draw[line width=3.5pt,-{Triangle[width=16pt,length=18pt]}]
    (500.7,113.1) ++(-144.6:30.01) arc[start angle=-144.6, end angle=-63.74, radius=30.01] -- (526.9,98.4);

  %% ════ 右框：x 处的小回环箭头（头朝下、指向圆点）════
  \draw[-{Triangle[width=16pt,length=18pt]}]
    (864.0,104.0) .. controls (864.0,88.0) and (874.0,70.0) .. (893.8,73.0)
                  .. controls (910.0,76.0) and (906.0,100.0) .. (884.0,112.0)
                  .. controls (879.0,116.0) and (876.0,117.0) .. (871.0,119.0);

  %% ════ 标签 ════
  %% 字母 a + 下标：**分开两个节点**，与 7.1 同一套写法（字母 30pt / 下标 22pt）
  %% 位置：字母在圆点左 5.3、下标在圆点右 7.1；字母 y=158.0、下标 y=168.0（实测）
  \foreach \dx/\sub in {108.8/1, 204.6/2, 306.1/3, 403.0/4, 472.0/5}{
    \node[anchor=center,inner sep=0pt,fill=white,font=\fontsize{30}{30}\selectfont]
      at (\dx-5.3,158.0) {\textsf{\textbf{\textit{a}}}};
    \node[anchor=center,inner sep=0pt,fill=white,font=\fontsize{22}{22}\selectfont]
      at (\dx+7.1,168.0) {\textsf{\textbf{\sub}}};
  }

  %% 映射 g —— 原稿写成了「9」，两处
  %% y 从实测中心 64.8 抬到 62：白底方框否则会切到弧顶（弧顶 y=75.9）
  \node[anchor=center,inner sep=1.5pt,fill=white,font=\fontsize{32}{38}\selectfont]
    at (328.3,62.0) {\textsf{\textbf{\textit{g}}}};
  \node[anchor=center,inner sep=1.5pt,fill=white,font=\fontsize{32}{38}\selectfont]
    at (924.5,83.4) {\textsf{\textbf{\textit{g}}}};

  %% 右框里的点标签 x
  \node[anchor=center,inner sep=1.5pt,fill=white,font=\fontsize{30}{30}\selectfont]
    at (850.4,157.1) {\textsf{\textbf{\textit{x}}}};

  %% 框标签 —— 原稿把左框的 B 拆成两段曲线，渲染出来像 θ；这里一律用文字
  \node[anchor=center,inner sep=1.5pt,font=\fontsize{30}{30}\selectfont]
    at (360.2,253.8) {\textsf{\textbf{\textit{B}}}};
  \node[anchor=center,inner sep=0pt,font=\fontsize{30}{30}\selectfont]
    at (833.7,253.3) {\textsf{\textbf{\textit{A}}}};
  \node[anchor=center,inner sep=0pt,font=\fontsize{30}{30}\selectfont]
    at (861.1,255.7) {{\mfallback\textbf{−}}};
  \node[anchor=center,inner sep=0pt,font=\fontsize{30}{30}\selectfont]
    at (888.4,253.3) {\textsf{\textbf{\textit{B}}}};

  % 钉死画布：否则超出的图元/控制点会把页面撑大，1:1 回比就失准
  \pgfresetboundingbox
  \path[use as bounding box] (0,0) rectangle (1037,278);
\end{tikzpicture}
\end{document}
